\section{구성요소 --- 하드웨어 및 소프트웨어 요소}

본 장은 ARES 시스템을 이루는 하드웨어 요소와 소프트웨어 요소를 조감한다.
세부 분석은 3장(하드웨어 핀)과 4장(펌웨어), 5\textasciitilde 8장(웹 계층)에서
이어진다.

\subsection{시스템 계층 구조}

ARES는 ``웹 클라이언트 --- 무선 링크 --- 임베디드 펌웨어 --- 물리 하드웨어''의
4계층으로 볼 수 있다.

\begin{longtable}{L{3.0cm} L{5.2cm} L{6.3cm}}
\toprule
\rowcolor{tablehead}\textbf{계층} & \textbf{구성요소} & \textbf{핵심 기술} \\
\midrule
\endhead
프레젠테이션 & 블록 편집기, 대시보드, 3D 랜딩/시뮬레이터 & HTML5, CSS, Google Blockly, three.js \\
애플리케이션 & 명령 생성기, 상태 관리, 로거, AI 도우미 & ES 모듈 JavaScript \\
통신 & Web Bluetooth $\leftrightarrow$ BLE UART & GATT UART 서비스(0xFFE0/0xFFE1), 20바이트 청크 \\
펌웨어 & 명령 처리기, 하드웨어 추상화, 설정 & MicroPython, 싱글턴 패턴 \\
물리 & 모터/센서/표시/음향/발사 & RP2040 GPIO, PWM, I2C, ADC \\
\bottomrule
\end{longtable}

각 계층은 바로 아래 계층에만 의존하며, 인접하지 않은 계층과 직접 결합하지
않는다. 예컨대 프레젠테이션 계층의 블록은 통신 계층의 BLE 세부를 알지 못하고,
통신 계층은 펌웨어 핸들러의 구현을 알지 못한다. 이 단방향 의존(그림~\ref{fig:layers})은
계층별 교체\textbf{·}시험을 쉽게 한다. 특히 통신 계층은 ``전송 가능한 텍스트 명령''만
다루므로, 실기(BLE) 대신 시뮬레이터(3D)로 대상을 바꿔도 상위 계층이 동일하게
동작한다(9장 시뮬레이션 분기 참조).

\begin{diagram}{4계층 단방향 의존 구조}{fig:layers}
\node[dom, text width=95mm] (l1) {프레젠테이션 — 블록 편집기 · 대시보드 · 3D 시뮬레이터};
\node[dom, text width=95mm, below=6mm of l1] (l2) {애플리케이션 — 명령 생성 · 상태 관리 · 로거 · AI 도우미};
\node[dom, text width=95mm, below=6mm of l2] (l3) {통신 — Web Bluetooth $\leftrightarrow$ BLE UART (0xFFE0/0xFFE1, 20B 청크)};
\node[dom, text width=95mm, below=6mm of l3] (l4) {펌웨어 — CommandProcessor · 하드웨어 추상화 · 설정};
\node[hw, text width=95mm, below=6mm of l4] (l5) {물리 — RP2040 GPIO · PWM · I2C · ADC};
\foreach \a/\b in {l1/l2, l2/l3, l3/l4, l4/l5}
  \draw[dep] (\a) -- (\b);
\node[font=\scriptsize\sffamily\color{accentcopper}, right=2mm of l3.east, anchor=west, text width=18mm]
  {텍스트 명령 계약};
\end{diagram}

\subsection{하드웨어 요소}

로버에 탑재되는 하드웨어 모듈은 모두 \emph{선택적}이며, 펌웨어 설정 파일
\code{system.json}을 통해 개별적으로 활성화\textbf{·}비활성화할 수 있다.

\begin{longtable}{L{2.6cm} L{2.6cm} L{8.8cm}}
\toprule
\rowcolor{tablehead}\textbf{모듈} & \textbf{드라이버 클래스} & \textbf{역할} \\
\midrule
\endhead
서보 휠 & \code{KSWheel} & 좌\textbf{·}우 연속회전 서보로 전진/후진/회전 (차동 구동) \\
DC 모터 & \code{KSDCMotor} & 메인 구동 모터, PWM H-브리지 속도 제어 \\
부저 & \code{KSBuzzer} & PWM 톤 생성, 비차단(non-blocking) 재생, 멜로디 지원 \\
LED 배열 & \code{KSLeds} & 6개 LED, PWM 밝기 제어(0.0\textasciitilde 1.0), 스와이프 효과 \\
거리 센서 & \code{KSDistance} & HC-SR04 초음파, cm 단위 거리 측정 \\
자기 센서 & \code{KSMagSensor} & 자기장 검출(디지털, 풀업) \\
발사기 & \code{KSGun} & BB탄 발사 모터, 소프트 스타트 \\
OLED & \code{SSD1306\_I2C} & 128$\times$64 I2C 디스플레이, 텍스트/아이콘 출력 \\
\bottomrule
\end{longtable}

\subsection{소프트웨어 요소}

\subsubsection{웹 클라이언트 (\code{Web/})}

ES 모듈로 분리된 JavaScript 파일들이 명확한 책임 경계를 가진다.

\begin{longtable}{L{3.4cm} L{11.0cm}}
\toprule
\rowcolor{tablehead}\textbf{파일} & \textbf{역할} \\
\midrule
\endhead
\code{main.js} & 앱 초기화, 라우팅, Blockly 워크스페이스 구성, 차시/미션 로딩 \\
\code{blocklyconfig.js} & 커스텀 블록 정의 및 카테고리(툴박스) \\
\code{commandexecutor.js} & 블록 $\to$ 명령 문자열 변환 및 전송, 제어흐름 처리 \\
\code{bluetooth.js} & \code{BluetoothManager}: BLE 연결, 청크 송수신, 응답 처리 \\
\code{constants.js} & UUID, 타이밍 상수, 기본 시스템 설정 \\
\code{state.js} & 전역 상태(연결, 변수, 실행 플래그) \\
\code{elements.js} & DOM 요소 참조 캐시 \\
\code{logger.js} & 통신 로그 UI \\
\code{simulation.js} & three.js 기반 3D 시뮬레이터 \\
\code{ai\_helper.js}, \code{ai\_helper\_grammar.js} & 자연어 $\to$ 블록 변환(문법 기반) \\
\code{romanize.js} & 한글 $\to$ 로마자 변환(OLED 텍스트용) \\
\code{mobile-preview.js} & 모바일 프레임 미리보기 \\
\bottomrule
\end{longtable}

\subsubsection{피코 펌웨어 (\code{Pico/})}

\begin{longtable}{L{3.4cm} L{11.0cm}}
\toprule
\rowcolor{tablehead}\textbf{파일} & \textbf{역할} \\
\midrule
\endhead
\code{main.py} & 부팅 시퀀스, UART 수신 메인 루프, 응답 전송 \\
\code{process\_data.py} & \code{CommandProcessor}: 명령 파싱 및 하드웨어 분배 \\
\code{hardware.py} & \code{RobotHardware} 싱글턴: 모듈 초기화/수명 관리 \\
\code{system\_config.py} & \code{SystemConfig}: \code{system.json} 입출력, 설정 관리 \\
\code{pins.py} & GPIO 핀 중앙 정의 \\
\code{wheel.py}, \code{dcmotor.py}, \code{buzzer.py}, \code{leds.py}, \code{gun.py} & 구동/출력 드라이버 \\
\code{ultrasonic.py}, \code{magsensor.py} & 센서 드라이버 \\
\code{ssd1306.py}, \code{icon.py} & OLED 드라이버 및 아이콘 비트맵 \\
\code{system.json} & 모듈 활성 플래그, 핀 할당, 보정값 저장 \\
\bottomrule
\end{longtable}

\subsubsection{외부 라이브러리 및 자산}

\begin{itemize}
  \item \textbf{Google Blockly} (v11) --- 시각적 블록 코딩 엔진.
        오프라인 빌드 시 \code{vendor/}에 로컬 사본으로 동봉.
  \item \textbf{three.js}\footnote{\textbf{three.js}: 웹 브라우저에서 WebGL을 쉽게 다루도록
        해주는 자바스크립트 3D 그래픽 라이브러리.} --- WebGL\footnote{\textbf{WebGL(Web
        Graphics Library)}: 웹 브라우저에서 별도 플러그인 없이 GPU 가속으로 3D\textbf{·}2D
        그래픽을 렌더링하는 표준 API.} 3D 렌더링.
        \code{vendor/three-bundle.min.js}로 사전 번들(로컬, CDN 비의존).
  \item \textbf{meshopt decoder}\footnote{\textbf{meshopt}: 3D 메시(정점\textbf{·}면) 데이터를
        작게 압축하고 빠르게 복원하는 기법\textbf{·}라이브러리. 모델 파일 용량을 줄여 로딩을
        앞당긴다.} --- \code{EXT\_meshopt\_compression}으로 압축된
        GLB\footnote{\textbf{GLB}: glTF 3D 모델을 하나의 바이너리 파일로 담은 포맷. 형상\textbf{·}
        재질\textbf{·}애니메이션을 함께 포함한다.} 모델 디코딩.
  \item \textbf{GLB 3D 모델} (\code{Web/Mesh/}) --- 로버, 알비, 신호등, 발사대 등
        미션 환경 모델.
\end{itemize}

\subsection{모듈 활성화 패턴}

모든 하드웨어 모듈은 동일한 안전 초기화 패턴을 따른다. 이는 일부 부품이
없거나 고장 난 상태에서도 시스템 전체가 부팅되도록 보장한다.

\captionof{listing}{hardware.py의 모듈 초기화 패턴}
\begin{minted}{python}
if sys_config.is_module_enabled('wheel'):
    try:
        from wheel import KSWheel
        self.wheel = KSWheel()
    except Exception as e:
        print("Wheel init failed:", e)
        self.wheel = None   # 실패 시 None -> 사용 측에서 건너뜀
\end{minted}

명령 처리 측(\code{process\_data.py})은 사용 전 \code{if not robot.wheel: return 0}와
같이 모듈 존재 여부를 확인하므로, 비활성 모듈에 대한 명령은 조용히 무시된다.
이 ``설정 검사 $\to$ 안전 초기화 $\to$ 사용 측 널 검사''의 3단 방어
(그림~\ref{fig:modinit})가 부분 고장\textbf{·}부분 구성 상황에서도 시스템 부팅과
운용을 보장하는 핵심 패턴이다.

\begin{diagram}{모듈 활성화 \textbf{·} 안전 초기화 연계 흐름}{fig:modinit}
\node[store, text width=20mm] (json) {system.json};
\node[boxw, right=12mm of json] (check) {is\_module\_\\enabled()};
\node[boxw, right=12mm of check] (try) {try: import\\+ 인스턴스화};
\node[hw, right=12mm of try, text width=20mm] (ok) {self.\textsl{mod}\\= 객체};
\node[boxw, below=8mm of try] (fail) {except:\\self.\textsl{mod}=None};
\node[boxw, right=12mm of ok, text width=24mm] (use) {process\_data:\\if not robot.\textsl{mod}: return 0};

\draw[flow] (json) -- node[lbl,above]{플래그} (check);
\draw[flow] (check) -- node[lbl,above]{ON} (try);
\draw[flow] (try) -- node[lbl,above]{성공} (ok);
\draw[flow] (try) -- node[lbl,right]{실패} (fail);
\draw[flow] (ok) -- (use);
\draw[back] (fail.east) to[bend right=20] node[lbl,below]{None $\Rightarrow$ 무시} (use.south west);
\end{diagram}

\begin{teachbox}{선택적 모듈과 안전 초기화}
\teachhead{설계 의도}
\begin{itemize}
  \item 모든 하드웨어를 \emph{선택적}으로 만들어, 부품 구성이 다른 로버도 같은 펌웨어로 돌린다.
  \item ``설정 검사 $\to$ try/except 초기화 $\to$ 사용 측 None 검사''의 3단 방어로 일부 부품이 없거나 고장 나도 전체가 부팅된다.
\end{itemize}
\teachhead{분석할 때 핵심 질문}
\begin{itemize}
  \item 어떤 모듈이 빠졌을 때 시스템은 멈추는가, 무시하는가? 왜 그렇게 했는가?
  \item 이 3단 방어가 없다면 어떤 문제가 생길까?
\end{itemize}
\teachhead{점검 요소}
\begin{itemize}
  \item \code{system.json}에서 한 모듈을 꺼 보고, 관련 명령이 어떻게 조용히 무시되는지 직렬 로그로 확인하라.
\end{itemize}
\end{teachbox}

